(18). This divergence, which is present even for standing waves ($r = \infty$), has its origin
in the Regge-Wheeler choice of gauge. It is a coordinate-dependent divergence which
can be removed by an infinitesimal coordinate transformation (change of gauge) and,
consequently, it is \textit{a divergence with no physical reality}. We have verified that it has no
physical reality by calculating---using a computer program of Thorne and Zimmerman
(1967)---the Riemann curvature invariants for the metric perturbation (18). We find
that, aside from exponential divergences associated with the damping time, $\tau$, the
perturbation in the curvature invariants has the acceptable form
\begin{equation}
  \delta R_I \sim 1/r \;.
\end{equation}

Edelstein (1967), in studying non-radial perturbations of the Schwarzschild geometry,
has independently verified that the divergence of order $r$ in $\delta g_{\mu\nu}/g_{\mu\nu}$ is non-physical. His
approach was to exhibit explicitly a coordinate transformation (change of gauge)
which made the amplitude of the exterior eigenfunctions go to zero for large $r$. We
shall not have need of this change of gauge in the present paper, but it is comforting
to know of its existence.

Summarizing the results of \S~III: To each choice of the complex eigenfrequency,
$\omega_n$, there is a unique eigenfunction which is physically acceptable at both $r = 0$ and
the surface of the star. That eigenfunction has the form (18) far from the star, which
is always a physically acceptable form if one allows for the possibility of both incoming
and outgoing gravitational waves.

\section{THE REAL PULSATIONS OF A RELATIVISTIC STELLAR MODEL}

The physical pulsations of a relativistic stellar model are \textit{real} linear combinations of
the complex normal modes discussed above and of their complex conjugates. We shall
describe these physical pulsations, first in words, and then in terms of the mathematics
of complex normal modes.

When an equilibrium configuration is perturbed, all of the perturbation energy is
concentrated initially in the motion of the fluid; there is no gravitational radiation.
However, as time passes the pulsation energy is gradually converted to gravitational
waves and radiated away. At any particular moment of time during the radiating phase,
the external gravitational waves have an amplitude which increases with radius until the
wave front is reached and which vanishes (no gravitational radiation) beyond the wave
front. The wave front moves forward with the speed of light carrying information about
the existence of the perturbation with it.

The mathematical description which accompanies these words is conveniently divided
into two regions: The region far behind the wave front both in distance and time, where
a very simple mathematical description suffices; and the region near the wave front,
where a more complicated description is necessary.

\subsection*{a) \textit{The Region Far behind the Wave Front}}

Far behind the wave front the pulsation and the gravitational waves consist of a
linear combination of discrete, complex normal modes with purely outgoing waves
(``\textit{outgoing normal modes}''):
\begin{equation}
  K(r,t) = \sum_n A_n \bigl[ K_n^{(0)}(r)\, e^{i\omega_n t} + K_n^{(0)*}(r)\, e^{-i\omega_n^* t} \bigr]
  \quad \text{and similarly for } H_0, W, V \;.
\label{eq:21}
\end{equation}

Here $K_n^{(0)}$ is the eigenfunction for the $n$th outgoing normal mode, $\omega_n$ is the
corresponding eigenfrequency, $A_n$ is a real amplitude, and the sum is over complex
conjugate pairs. Note that the pulsations described by equation (21) are far from the
most general ones possible: They are restricted to a particular spherical harmonic
$l$, $M \neq 0$, $\tau =$
